\subsection{Closed-loop agent architecture}

AdsMind is implemented as a deterministic state machine with four explicit roles: a planner, a validator, a tool executor, and a final analyzer. The planner proposes a structured adsorption hypothesis in JSON format, including the canonical site type (\textit{ontop}, \textit{bridge}, or \textit{hollow}), the intended surface symbols, and the adsorbate binding indices. The validator then enforces chemistry- and schema-level constraints before any expensive physics call is made. In the current implementation, this validation step rejects duplicate plans, invalid binding-index choices, and site/index combinations that are inconsistent with the visible heavy-atom representation of the adsorbate.

The tool executor instantiates the accepted plan on the slab, performs physical relaxation, and returns a structured post-relaxation analysis rather than only a scalar energy. The analyzer records whether the final structure remained bound, whether dissociation or rearrangement occurred, and whether the realized coordination environment differs from the planner's intent. We refer to this latter event as a \textit{Chemical Slip}: the agent asked for one adsorption environment, but relaxation converged to another. The final analyzer updates the running memory with this outcome and either requests another hypothesis or terminates the episode.

\subsection{Prompting, memory, and ablation switches}

The planner prompt is assembled from conditional sections instead of a monolithic fixed template. Three runtime switches control the agentic feedback exposed to the language model. \textit{Slip feedback} determines whether prior slips are framed as explicit evidence that the originally planned site is unstable. \textit{FORBID} determines whether slipped or otherwise failed site intents are converted into an explicit negative constraint in the planner history. \textit{Termination} determines whether convergence-to-known-best signals are surfaced to the planner and whether the planner is allowed to stop early after repeated convergence.

These switches define the ablation variants used in the Results section. The \textit{Full} setting enables all three mechanisms. \textit{$-$Slip} removes slip-conditioned guidance and the associated negative constraint. \textit{$-$Forbid} keeps slip detection but suppresses the explicit prohibition against retrying a failed site family. \textit{$-$Term} disables convergence-based early stopping while preserving the rest of the loop. \textit{1-Shot} disables all three mechanisms and sets the maximum attempt budget to one.

\subsection{Semantic-to-geometric execution}

Each accepted plan is translated into an atomistic starting structure through a surrogate-SMILES geometry layer. This layer maps the selected binding atom(s) to proxy adsorption markers, constructs physically valid initial placements for \textit{ontop}, \textit{bridge}, and \textit{hollow} adsorption, and samples a capped set of candidate surface sites and conformers. In the benchmark configuration used here, the default cap is eight sampled surface coordinates per semantic site family and four conformers per site, with one top-ranked candidate carried into full relaxation.

The physical stage uses a two-step evaluation procedure. Candidate structures are first screened cheaply to eliminate obviously unstable configurations and to rank promising candidates. The best surviving candidate is then fully relaxed and analyzed. The post-relaxation analysis reconstructs the actual local coordination environment, detects dissociation and bond changes, and writes structured artifacts that can be consumed directly by the next planning step or by offline benchmark tooling.

\subsection{Benchmark and ablation protocol}

All agent-side benchmark runs in this manuscript use frozen experiment configurations rather than mutable product defaults. Three reasoning backends are evaluated: Gemini~2.5~Pro (\texttt{gemini-2.5-pro}), Grok-4 (\texttt{grok-4-0709}), and GPT-5.4 (\texttt{gpt-5.4-2026-03-05}), all run at zero temperature through each provider's official API endpoint. For the local Apple Silicon benchmark, the physical backend is locked to a conservative MACE configuration: \texttt{model=small}, \texttt{precision=float32}, \texttt{device=cpu}, and dispersion disabled. This choice reflects the practical hardware constraints of the benchmark environment and is held fixed across all reported ablations and backend comparisons. Each frozen configuration records the exact model identifier, API transport variant, and key resolution path so that provenance is traceable per run.

The representative ablation matrix uses five locked cases (01, 02, 09, 14, and 19), chosen to span simple atomic adsorption, multi-atom intermediates, hydroxyl adsorption, and a difficult large-adsorbate case. Full runs use a maximum of five attempts. The 1-shot baseline uses the same tooling and physical backend but restricts the planner to a single attempt, allowing us to isolate the value of closed-loop refinement from the value of the base planner alone. All batch experiments are executed sequentially, with fixed seeds and structured result serialization for reproducibility.
